% chktex-file 44 32 36 38 18 11 8 26


\medskip

\exods{0}
%AN1J2 exo3

\emph{Cet exercice porte sur les bases de données relationnelles et le langage SQL.}

L’énoncé de cet exercice utilise les mots du langage SQL suivant :

\verb|DELETE, FROM, INSERT INTO, JOIN, SELECT, SET, UPDATE, VALUES, WHERE.|

\par 
Les clés primaires seront soulignées et les clés étrangères seront précédées d’un \verb|#|.\par

Le satellite GAIA a pour mission de cartographier un très grand nombre d’objets
autour du Système Solaire. Régulièrement un catalogue est produit pour publier les
données obtenues. Il est disponible sous différents formats dont par exemple sous forme
de fichier \verb|csv|. \verb|Num_Objet| identifie chaque objet cartographié de manière unique.

\begin{table}[h]
\centering
\includegraphics[width = 0.9\linewidth]{donnees_astro.png}
\caption{Données Astronomiques sur les Systèmes Stellaires Proches}
\label{tab:astronomie}
\end{table}

Pour manipuler plus facilement les données, un chercheur utilise un système de base de données relationnelle, dans lequel il crée le schéma relationnel de la table Gaia :

\begin{verbatim}
Gaia (
    Num_Objet : Int, Num_Systeme : Int, Nom_Systeme : String, Type_Objet : String, 
    Nom_Objet : String, Ascension_Droite : Real, Declinaison : Real, Parallaxe : Real, 
    Nom_SIMBAD : String)
\end{verbatim}

\medskip

\section*{PARTIE A}

\medskip

\begin{enumerate}[resume=monDS]
    \item Justifier que l’attribut \verb|Num_Objet| peut être choisi comme clé primaire de la table \verb|Gaia|.
\end{enumerate}

Le type de l’objet ( \verb|Type_Objet| ) n’est pas une information directement compréhensible et le chercheur décide de créer une nouvelle table appelée \verb|Type|. La clé primaire est \verb|Type_Objet|.\par

Soit la table \verb|Type| suivante :

\begin{table}[h]
\centering
\includegraphics[width = 0.4\linewidth]{donnees2.png}
\caption{Table Type}
\label{tab:type}
\end{table}

\begin{enumerate}[resume=monDS]
    \item Proposer le schéma relationnel de la table \verb|Type| en soulignant la clé primaire.
\end{enumerate}

Dans ce qui suit, la table \verb|Gaia| a été modifiée de manière à ce que le champ \verb|Type_Objet| soit une clé étrangère associée à la clé primaire de même nom de la table \verb|Type|.


\medskip
\thispagestyle{DS_LP}

\section*{PARTIE B}

\medskip

\begin{enumerate}[resume=monDS]
    \item Expliquer pourquoi le code SQL suivant ne fonctionne pas :
    
    \begin{verbatim}
INSERT INTO Type VALUES ('BD', 'Trou Noir') ;
    \end{verbatim}

    \item Indiquer le résultat de la requête suivante exécutée sur l’extrait présenté :
    
    \begin{verbatim}
SELECT Nom_Objet, Parallaxe FROM Gaia WHERE Type_Objet = 'Planet';    \end{verbatim}


\item Écrire la requête qui permet de récupérer le nom du système, le nom de l’objet et le libellé du type pour des objets ayant une parallaxe supérieure à 400 millisecondes d’arc et étant des \verb|'Étoile'|.
\item Le chercheur a besoin du nom, du libellé de leur type, de la déclinaison et de la parallaxe de tous les objets du système stellaire Alpha du Centaure ( \verb|'alf Cen'| ). Écrire la requête qui permet d’extraire ces informations de la base de données.

\end{enumerate}
